\item En un cilindro, una muestra de un gas ideal con número de moles $n$ experimenta un proceso adiabático. 
(a) Partiendo de la expresión $W = -\int P dV$ y empleando la condición $PV^\gamma = \text{constante}$, demuestre que el trabajo realizado sobre el gas es
$$ W = \left( \frac{1}{\gamma - 1} \right) (P_fV_f - P_iV_i) $$
(b) Partiendo de la primera ley de la termodinámica, demuestre que el trabajo efectuado sobre el gas es igual a $nC_V(T_f - T_i)$. 
(c) ¿Estos dos resultados son consistentes entre sí? Explique.
